% RESEARCH_DRAFT_SHA256: b9d06c326bbbef5be0e53b3e18ce05d53580c83cdd91195c45337e17ce667c25
% Layout and source-reference wrappers expanded for the web renderer.
\documentclass[11pt]{article}
\usepackage[margin=25mm]{geometry}
\usepackage{fontspec,unicode-math}
\setmainfont{Latin Modern Roman}
\setmathfont{Latin Modern Math}
\usepackage{amsmath,amssymb,mathtools}
\usepackage{xurl,hyperref}
\hypersetup{colorlinks=true,urlcolor=blue,linkcolor=blue}
\setcounter{secnumdepth}{0}
\setlength{\parindent}{0pt}
\setlength{\parskip}{5pt}
\setlength{\emergencystretch}{3em}
\newcommand{\sref}[2]{\hyperlink{source:#1:#2}{[S#2]}}
\newcommand{\cataloguescope}[1]{\paragraph{Source question.}#1}
\newcommand{\floor}[1]{\left\lfloor#1\right\rfloor}
\begin{document}
\section*{Research-attempt provenance}
This standalone research attempt retains the complete original definition body below,
including its historical source and preparation notes. Those notes describe the
earlier catalogue preparation; the new attempt and its verification are identified
separately after that body. No original ID, catalogue status, or source statement is changed.

% UnsolvedMath ID 30003185; source catalogue code OWR-14745-003
\section{Complexity of Bray’s Involution-Centralizer Algorithm}\label{30003185}
{\small UnsolvedMath ID 30003185 · Computer Science}
\subsection{Definitions and mathematical statement}

\paragraph{Shared notation.}Standard number systems are denoted by $\mathbb N,\mathbb Z,\mathbb Q,\mathbb R,\mathbb C$; index ranges, computational models and specialized notation follow the cited problem source and the definitions below.


A strong involution $t$ has order two and both its $+1$ and $-1$ eigenspaces have dimensions between $n/3$ and $2n/3$. Bray\textquotesingle s method samples independent uniform $g\in G$, sets $y=tt^g$ with $t^g=g^{-1}tg$, and extracts centralizer elements from the odd/even order cases. The centralizer is $C_G(t)=\{h:ht=th\}$.

The group $\mathrm{GL}_n(\mathbb Z)$ consists of integer matrices with integer inverses, and $\mathrm{SL}_n$ imposes determinant one. The subgroup generated by matrices consists of finite products of them and their inverses. A centralizer consists of elements commuting with the specified element.
\paragraph{Statement recorded in the supplied source.}
For a finite classical group of dimension $n$ over a field of odd characteristic, does Bray's algorithm for constructing the centralizer of a strong involution succeed with high probability after examining only $O(\log n)$ random elements? \href{https://doi.org/10.4171/owr/2016/37}{[S2]}
\paragraph{Scope qualification.}
The original report treats the even-case construction of a subgroup containing the second derived subgroup $C_G(t)^{\prime\prime}$ as sufficient for its modified aim. This is not silently strengthened to generating every element of the full centralizer.
\paragraph{Source status.}
The supplied catalogue labels this record partially solved. Its source-status metadata was checked against the supplied snapshot on 2026-10-11. This import does not independently establish current open status or mathematical correctness. The catalogue literature assessment is dated 2026-08-21.
\subsection{Short English statement}
Prove logarithmic random-sampling complexity for the stated strong-involution centralizer algorithm.
\subsection{Sources}
\begin{itemize}
\item[S1] Ulam AI and the UnsolvedMath Contributors, UnsolvedMath: A Curated Collection of Open Mathematics Problems, record 30003185 (OWR-14745-003). Catalogue curation and metadata: CC BY 4.0. The supplied source question is retained; definitions and the English summary are editorial additions.
\url{https://huggingface.co/datasets/ulamai/UnsolvedMath}.
\item[S2] Computational Group Theory (2017), Oberwolfach Reports, DOI 10.4171/owr/2016/37, pp. 2123--2169
\url{https://doi.org/10.4171/owr/2016/37}.
\item[S3] Dixon--Praeger--Seress
\url{https://www.sciencedirect.com/science/article/pii/S002186931730635X}.
\item[S4] Glasby--Praeger--Roney-Dougal
\url{https://doi.org/10.1016/j.jalgebra.2019.09.009}.
\item[S5] Original Oberwolfach Report 37/2016, Cheryl E. Praeger, pp. 2160--2161
\url{https://ems.press/content/serial-article-files/46642}.
\textit{Underlying Oberwolfach Reports material retains its source-specific licence, including CC BY-SA 4.0 where applicable; see THIRD\_PARTY\_NOTICES.md.}
\end{itemize}
\paragraph{Preparation and review.}
Original question attributed to the cited source and UnsolvedMath contributors. Codex (GPT-6; exact serving revision not exposed) assisted with definitions, English summaries, source inspection and formatting. No new human mathematical authorship or independent proof verification is asserted. Canonical ID 30003185 and source code OWR-14745-003 checked against the complete supplied registry and existing repository statements; known flagged repetitions and obvious identical conjectures excluded. Full mathematical-equivalence/literature review has not been performed. Statement preparation and final review followed contribution prompts C and D; automated checks verify structure and identity, not open status.
\label{end:30003185}

\section{Literature and context check}
OWR 37/2016, pp.~2160--2161, was read directly. It already proposes a lower bound $c'/\log n$ for obtaining a useful even-case involution on each eigenspace, reports the special linear case, and distinguishes generating a subgroup containing $C_G(t)''$ from generating the full centralizer. Its cited generation theorems and later journal papers have not been independently rederived or freshly checked here. This attempt audits the sampling implication, not those group-specific probability estimates.
\section{Research attempt: Sampling amplification and its missing uniformity}

\paragraph{A quantified conditional sampling lemma.}
Let $n\ge2$. Suppose each independent uniform input $g$ has probability at least $p_n>0$ of producing a useful even-case output for each specified sign $\varepsilon\in\{+,-\}$. The two signs within a trial need not be independent. Fix a positive integer $c$ and $0<\delta<1$. After $N$ trials let $X_\varepsilon$ count useful outputs for sign $\varepsilon$. Independence across trials gives, with $a=1-e^{-1}>1/2$,
\[
 \mathbb E e^{-X_\varepsilon}\le (1-a p_n)^N\le e^{-aNp_n},
 \qquad
 \Pr(X_\varepsilon<c)\le e^{c-aNp_n}.
\]
The first formula follows by multiplying $\mathbb E e^{-I}=1-a\Pr(I=1)$ for the Bernoulli indicators; the second is Markov's inequality applied to $e^{-X_\varepsilon}$, using the weaker threshold $e^{-c}$. Hence
\[
 N\ge \frac{2}{p_n}\bigl(c+\log(2/\delta)\bigr)
 \quad\Longrightarrow\quad
 \Pr(X_+<c\text{ or }X_-<c)\le\delta.
\]
The union bound, not independence of the two counts, supplies the last step. If $p_n\ge a_0/\log n$ for a fixed $a_0>0$, the sufficient sample count is
\[
 O\bigl(\log n\,[c+\log(1/\delta)]\bigr).
\]
For fixed $c$ and fixed error tolerance this is $O(\log n)$. If a further group-generation assertion fails with conditional probability at most $\eta$ after obtaining the required outputs, the total failure is at most $\delta+\eta$; that assertion has not been supplied by the counting lemma.

\paragraph{The attempted stronger inference fails.}
A probability lower bound of order $1/\log n$ alone does not imply success tending to one from $C\log n$ trials with fixed $C$. Indeed the independent Bernoulli model with exact success probability $a_0/\log n$ (for sufficiently large $n$) has
\[
 \Pr(X=0)=\left(1-\frac{a_0}{\log n}\right)^{\lfloor C\log n\rfloor}
 \longrightarrow e^{-a_0C}>0.
\]
For completeness, $\log(1-z)/z\to-1$ follows by integrating $1/(1-z)$ and bounding it between $1$ and $1/(1-z)$; this proves the displayed limit. Taking $\delta=n^{-b}$ in the sufficient bound yields $O((\log n)^2)$ for fixed $b>0$. This is an obstruction to deriving an asymptotically vanishing error from that lower bound alone, not a lower bound for Bray's actual distribution, which might be better.

\paragraph{Why qualitative strength is also insufficient.}
In $\mathrm{GL}_4(\mathbb F_q)$ with odd $q$, the diagonal involutions having two $-1$ and two $+1$ entries are strong. Any collection consisting only of these diagonal elements generates an abelian group, whereas, for example, the elementary matrices $I+E_{12}$ and $I+E_{23}$ do not commute: their products differ by $E_{13}$. Thus strong eigenspace dimensions alone cannot imply generation of a nonabelian target. Uniform conjugacy-class marginals alone are also insufficient if independence is dropped: choose one uniform conjugate $z$ of an involution and repeat $z$ in every position; every marginal is correct while the generated group has order at most two. Bray's independent inputs avoid this artificial dependence, but the needed distribution and generation hypotheses still require proof.
\section{Verification and adversarial audit}
The exponential-moment calculation was derived rather than quoted. The argument permits useful outputs on both signs in the same trial, checks $n\ge2$, and distinguishes fixed-error sampling from probability tending to one. The diagonal example is an auxiliary failure of an insufficient condition; it is not claimed to be the distribution of Bray outputs or a counterexample to the original conjecture.
\paragraph{Audit corrections.} Removed the invalid implication that $c'/\log n$ immediately gives asymptotic success probability one from a fixed multiple of $\log n$ trials. Kept the original source's possible fixed-tolerance interpretation open rather than declaring it erroneous.
\paragraph{External computation.} None; symbolic arguments and boundary checks only.
\section{Scope and remaining gaps}
No new uniform lower bound for useful even-case outputs in all odd-characteristic classical families has been proved, and no group-generation theorem has been supplied. The source's second-derived-subgroup aim remains distinct from the full centralizer. The first missing step is the actual family-uniform probability/distribution analysis, with an explicit meaning of high probability.
\paragraph{Integration review.} Codex (GPT-6; exact serving revision not exposed) reviewed the written arguments and expanded layout and source-reference wrappers. Repository metadata and the production web renderer were checked using Python 3.13 and JavaScript with MathJax 3.2.2. This model-assisted review is not independent human proof verification; no PDF compilation or formal proof checking is asserted.

\section{Completion Estimate}
\noindent\textbf{COMPLETION: 0\%}
The calculations below the preserved source statement investigate a concrete route,
but establish no substantiated new progress toward the entire original target.
Known special cases and failed sufficient conditions are not credited as a solution.
\section{References}
The original source references and catalogue attribution are preserved in the
Sources section above. The following URLs were actually consulted for this attempt;
retention of an additional original reference does not claim fresh consultation of it.
\begin{itemize}
\item \url{https://ems.press/content/serial-article-files/46642} (consulted 2026-10-11). 
\end{itemize}
\section{Model and contributor disclosure}
OpenAI gpt-6-astra, requested reasoning configuration ultra, generated this research
attempt on 2026-10-11. The exact serving revision is not exposed. Codex GPT-6
performed orchestration, repository-format review, and written-proof review of this attempt; its exact serving revision
is also not exposed. The mathematical route, boundary checks and final adversarial
self-audit followed contribution prompts A then B in
\texttt{docs/contribution\_prompts.md}, after reading \texttt{CONTRIBUTING.md}.
Python 3.14.2 assembled the standalone documents; any mathematical computation is
specified in Verification. No theorem prover was used. No human mathematical
authorship, substantial human proof correction, or independent human proof checking
is asserted. Model-assisted review, source retrieval, and structural checks do not
establish correctness or novelty. Full private reasoning transcripts are unavailable
as reproducibility artifacts; the complete arguments, inputs and outcomes actually
used are stated here or in the companion batch audit.
\section{Final status}
Unresolved as a new research contribution. This is a concrete route investigation,
not a claim of new complete resolution or a current-status certification. A known
answer reproduced for context does not become a new contribution. The catalogue's
source-status label is preserved; this attempt status does not assert that a
source-labelled solved question is open.
\end{document}
